\chapter*{Résumé}
\addcontentsline{toc}{chapter}{Résumé}

Avec l'évolution rapide des technologies d'Intelligence Artificielle, les modèles de langage ont prouvé une grande efficacité pour assimiler d'énormes quantités d'informations et les reformuler afin de répondre à des questions complexes. En revanche, les consultations juridiques traditionnelles et l'analyse de dossiers restent des processus très coûteux. Une utilisation appropriée de l'IA peut réduire ce coût de manière drastique pour le rendre presque symbolique.

Notre objectif à travers ce projet de fin d'année est de développer une plateforme numérique intégrée permettant aux utilisateurs d'obtenir des consultations juridiques précises, d'analyser leurs dossiers et de générer automatiquement des documents juridiques. Pour ce faire, nous avons d'abord mené une étude approfondie des besoins du marché et défini les fonctionnalités nécessaires pour résoudre les problématiques soulevées. Ensuite, nous avons modélisé l'ensemble du système à l'aide de diagrammes UML.

Nous avons été confrontés à un obstacle majeur : les modèles globaux (tels que Gemini, GPT, Claude) ne sont pas suffisamment entraînés sur le droit marocain, ce qui entraîne des « hallucinations » et des informations erronées. Pour y remédier, nous avons construit une architecture avancée basée sur la Génération Augmentée par la Recherche (RAG). Cela nous a permis de créer une base de données vectorielle avec ChromaDB, alimentée par la reconnaissance optique de caractères (OCR) via GPT-4o avec des optimisations spécifiquement adaptées aux documents marocains.

Sur le plan technique, nous avons utilisé FastAPI pour un backend performant, React pour les interfaces utilisateur, et PostgreSQL pour la gestion de la base de données relationnelle.

\noindent\rule[2pt]{\textwidth}{0.5pt}

{\textbf{Mots clés :}}
Consultations juridiques, Intelligence Artificielle, Hallucination, RAG, Base de données vectorielle, ChromaDB, GPT-4o, OCR, UML, FastAPI, React, PostgreSQL.
\\
\noindent\rule[2pt]{\textwidth}{0.5pt}